% Appendix: hosted engineering -- drafted by Claude (2026-10-21), revised after review.
\section{Calling the API during scoring: details}\label{app:engineering}

\paragraph{Probes.} A first network probe (Codabench 964074) was built to fail unless \texttt{api.anthropic.com} answered. It produced predictions for all 8 pairs at every budget, but Codabench later marked it Failed without giving a reason, so I do not rely on it. A second probe (968358) failed with a missing \texttt{anthropic} module, which showed that \texttt{requirements.txt} is not installed. A third (968649) used only the standard library and predicted one constant per budget that encoded what it found: whether a request with a dummy key got the expected 401 reply from the API, whether the submission folder is writable, and whether files written at one budget are visible at the next. All three checks passed at every budget.

\paragraph{Failure classes.} Each failed request is classified from its status code. Busy replies (429, 529) and requests that never left the machine are free: every worker pauses for the time the API asks for (at most 60 seconds), and the request is retried. Refusals of the key or account (401, 402, 403 and billing errors) write a marker file that stops new requests in every worker; replies already cached are still used. Requests the API can never answer (400, 404, 413, 422) leave the item unrated or the subject unguessed. Anything else counts as transient and may be billed, so each worker tries it at most twice.

\paragraph{Locks and the cap.} Before sending a request, a worker creates a lock file holding its process id and a random token, waits 50 ms and reads the file back; if two workers took over the same lock at once, only the one whose token survives sends the request. A lock older than 45 seconds, or one whose owner process has ended, is taken over. Each billed attempt writes a marker file first, and no request is sent once 1,000 markers exist. Requests run on daemon threads (8 per worker) with a 20-second socket timeout, so a hanging API never delays the end of a budget. Timed from my machine, a guess took 1.4 to 3.1 seconds and a rating 1.3 to 1.7 seconds, against the 4-second wait of a prediction.

\paragraph{The key.} The key never enters the repository; I add it to a private copy of the zip with a script that checks it first. It belongs to a workspace used for nothing else, with a \$50 monthly spending limit. The first version of that check accepted only keys beginning \texttt{sk-ant-api}, while mine begins \texttt{sk-ant-usr}, so the first V4 build would have run as V3 on Codabench; running the key script with my real key caught this before upload.

\paragraph{Tests.} The offline numbers of Section~\ref{sec:results} come from my predictor classes, refitted without the held-out benchmark. As a check of the packaging, I also ran the V4 and V5 submission folders through the harness, with their reply caches filled from the saved ratings and guesses so that they made no API calls. In the rated setting V5's folder scored $-0.0150$ (SE 0.0012) against V4's. Because these folders were fitted on all public benchmarks, the scores are in-sample; the size of the gain still agrees with the held-out test.

Before uploading V5 I ran 156 unit tests and the starting kit's zip checker and smoke test, on Windows and under Linux (also with four isolated workers). A rebuild under Linux gave the same SHA-256, and a harness run without a key matched V3 on all 576 pair results of the default setting. A live run on 6 public items, with the key held only in memory, gave valid guesses and the same predictions as the in-repository model, and wrote the key to no file. With a hanging API every prediction returned within about 4.1 seconds, and a simulation that killed four workers at every budget still gave valid predictions and one guess per subject and benchmark.
